\documentclass[../../main.tex]{subfiles}
\begin{document}

{\bf [解]}
\paragraph{(1)}

\begin{figure}[htb]
\centering
\begin{tikzpicture}[scale=2.5, >=stealth, every node/.style={font=\small}]

  % 1. 基準点 A の定義
  \coordinate (A) at (0, 0);

  % 2. 2本の半直線（接線上の点と回転後の直線 l 上の点）
  % 接線方向 (例: 傾き 50°)
  \coordinate (Pprime) at (50:1.4);
  % -30° 回転した直線 l 方向 (50° - 30° = 20°)
  \coordinate (P) at (20:1.5);

  % 3. 線分の描画
  % 接線方向
  \draw[thick, blue!80!black] (A) -- (Pprime);
  % 回転後の直線 l
  \draw[thick, red!80!black] (A) -- (P);

  % 4. 回転角 -30° のアーク
  \draw[->, thick, gray!80!black] (50:0.6) arc[start angle=50, end angle=20, radius=0.6];
  \node[right, font=\footnotesize] at (35:0.65) {$-30^\circ$};

  % 5. 頂点・点のプロットと縦ベクトル表記のラベル
  \fill (A) circle (0.8pt) node[below left] {$A$};
  
  % 接線上の点 (x', y')
  \fill[blue!80!black] (Pprime) circle (0.8pt);
  \node[above, font=\normalsize] at ($(Pprime) + (0, 0.08)$) {$\begin{pmatrix} x' \\ y' \end{pmatrix}$};

  % 直線 l 上の点 (x, y)
  \fill[red!80!black] (P) circle (0.8pt);
  \node[right, font=\normalsize] at ($(P) + (0.05, 0)$) {$\begin{pmatrix} x \\ y \end{pmatrix}$};

\end{tikzpicture}
\caption{点 $A$ を中心とする $-30^\circ$ の回転の様子}
\end{figure}

$A=(a,a^2)$における$y=x^2$の接線は$y=2a(x-a)+a^2$である．
この接線を$A$を中心に$-30^\circ$回転した直線が$l$である．$l$上の点$(x,y)$を$+30^\circ$回転すると接線上の点$(x',y')$に戻るから
\begin{align}
  \begin{pmatrix}
    x'-a\\
    y'-a^2
  \end{pmatrix}
=\begin{pmatrix}
   \cos\dfrac{\pi}{6} &-\sin\dfrac{\pi}{6}\\
   \sin\dfrac{\pi}{6} &\cos\dfrac{\pi}{6}
\end{pmatrix}
\begin{pmatrix}
  x-a\\
  y-a^2
\end{pmatrix}
=\begin{pmatrix}
   \dfrac{\sqrt{3}}{2} &-\dfrac{1}{2}\\
   \dfrac{1}{2}        & \dfrac{\sqrt{3}}{2}
\end{pmatrix}
\begin{pmatrix}
  x-a\\
  y-a^2
\end{pmatrix}
=\begin{pmatrix}
  \dfrac{\sqrt{3}}{2}(x-a)-\dfrac{1}{2}(y-a^2)\\[4pt]
  \dfrac{1}{2}(x-a)+\dfrac{\sqrt{3}}{2}(y-a^2)
\end{pmatrix}
\end{align}
であり，これが接線上の点であることから$y'-a^2=2a(x'-a)$，すなわち
\begin{align}
\dfrac{1}{2}(x-a)+\dfrac{\sqrt{3}}{2}(y-a^2)=2a\Bigl(\dfrac{\sqrt{3}}{2}(x-a)-\dfrac{1}{2}(y-a^2)\Bigr)
\end{align}
である．

整理すると，$l$の方程式は
\begin{align}
(\sqrt{3}+2a)y=(2\sqrt{3}a-1)x+a+2a^3-\sqrt{3}a^2 \label{eq:1}
\end{align}
である．


\paragraph{(2)}

\begin{figure}[htb]
\centering
\begin{tikzpicture}[scale=2.2, >=stealth, every node/.style={font=\small}]

  % --- 1. パラメータ設定と解析計算 ---
  % a の値 (例: a = 1.2)
  \def\a{1.2}
  % 直線 l と放物線 y=x^2 のもう一つの交点 B の x 座標
  % 直線 l の傾き m = (2\sqrt{3}a - 1) / (\sqrt{3} + 2a)
  \pgfmathsetmacro{\m}{(2.0*sqrt(3)*\a - 1.0) / (sqrt(3) + 2.0*\a)}
  % 解と係数の関係より x_B = m - a
  \pgfmathsetmacro{\xb}{\m - \a}
  \pgfmathsetmacro{\ya}{\a*\a}
  \pgfmathsetmacro{\yb}{\xb*\xb}

  % 座標定義
  \coordinate (O) at (0, 0);
  \coordinate (A) at (\a, \ya);
  \coordinate (B) at (\xb, \yb);
  \coordinate (C) at (\a, 0);

  % --- 2. 領域の塗りつぶし ---
  % S(a): x軸 (OC), 垂線 CA, 放物線 OA で囲まれる部分 (斜線ハッチング)
  \fill[pattern=north east lines, pattern color=black!40] 
    (O) -- (C) -- (A) 
    -- plot[domain=\a:0, samples=40] (\x, {\x*\x}) 
    -- cycle;

  % T(a): 線分 AB と放物線 AB で囲まれる部分 (薄い青色塗りつぶし)
  \fill[blue!20, opacity=0.8] 
    (B) -- (A) 
    -- plot[domain=\a:\xb, samples=40] (\x, {\x*\x}) 
    -- cycle;

  % --- 3. 補助線（点線） ---
  \draw[dashed, gray!70] (A) -- (C);
  \draw[dashed, gray!70] (B) -- (\xb, 0);

  % --- 4. 座標軸 ---
  \draw[->, thick] (-1.0, 0) -- (1.8, 0) node[right] {$x$};
  \draw[->, thick] (0, -0.2) -- (0, 2.0) node[above] {$y$};
  \node[below left] at (O) {$O$};

  % --- 5. 曲線・直線の描画 ---
  % 放物線 y = x^2
  \draw[very thick, blue!80!black, domain=-0.9:1.35, samples=60] 
    plot (\x, {\x*\x}) node[above right] {$y = x^2$};

  % 回転した直線 l
  \draw[very thick, red!80!black, domain=-0.75:1.35] 
    plot (\x, {\m*(\x - \a) + \ya}) node[right] {$l$};

  % --- 6. 軸上の目盛りラベル ---
  \node[below] at (\a, 0) {$a$};
  \node[below] at (\xb, 0) {$x_B$};

  % --- 7. プロット点とラベル ---
  \fill (O) circle (0.7pt);
  \fill (C) circle (0.7pt) node[below right] {$C(a,0)$};
  \fill[red!80!black] (A) circle (0.9pt) node[above right] {$A(a,a^2)$};
  \fill[red!80!black] (B) circle (0.9pt) node[above left] {$B$};

  % --- 8. 領域ラベル ---
  \node[fill=white, inner sep=1pt, font=\bfseries] at ({0.6*\a}, {0.25*\ya}) {$S(a)$};
  \node[fill=white, inner sep=1pt, font=\bfseries] at ({(\a+\xb)/2}, {(\ya+\yb)/2 - 0.25}) {$T(a)$};

\end{tikzpicture}
\caption{放物線 $y=x^2$ と直線 $l$ で囲まれる領域 $T(a)$ および領域 $S(a)$}
\end{figure}

\cref{eq:1}と$y=x^2$の交点の$x$座標は，$y=x^2$を\cref{eq:1}に代入した$x$の二次方程式
\begin{align}
(\sqrt{3}+2a)x^2 - (2\sqrt{3}a-1)x - a - 2a^3 + \sqrt{3}a^2 = 0
\end{align}
の解である．一方は$x=a$（点$A$）だから，解と係数の関係より，他方の解$\alpha$（点$B$の$x$座標）は
\begin{align}
  a\alpha&=\dfrac{-(a+2a^3-\sqrt{3}a^2)}{\sqrt{3}+2a}\\
\therefore\ 
  \alpha&=\dfrac{-2a^2+\sqrt{3}a-1}{2a+\sqrt{3}}\qquad(<a) \label{eq:2}
\end{align}
である．

従って，まず直線$l: y=\ell(x)$と$y=x^2$の囲む面積$T(a)$は
\begin{align}
 T(a) 
  &= \int_{\alpha}^{a} \left(y^2 - \ell(x)\right) dx \\
  &= \int_{\alpha}^{a} \left(y - \alpha \right)\left(y - a\right) dx \\
  &= \dfrac{1}{6}\left(a-\alpha\right)^3 \\
  &= \dfrac{1}{6}\left(a- \dfrac{-2a^2+\sqrt{3}a-1}{2a+\sqrt{3}} \right)^3 \quad (\because\text{\cref{eq:2}})\\
  &= \dfrac{1}{6}\left(\dfrac{4a^2+1}{2a+\sqrt{3}}\right)^3 \label{eq:3}
\end{align}
である．

一方，線分$OC$，$CA$と$y=x^2$で囲まれる面積は$x^2$の積分で
\begin{align}
S(a)=\int_0^ax^2dx=\dfrac{1}{3}a^3 \label{eq:4}
\end{align}
である．

よってこれらの比は
\begin{align}
\dfrac{T(a)}{S(a)}
  &=\dfrac{3}{6}\Bigl(\dfrac{4a^2+1}{2a+\sqrt{3}}\Bigr)^3\Big/a^3 \\
  &=\dfrac{1}{2}\Bigl(\dfrac{4a^2+1}{a(2a+\sqrt{3})}\Bigr)^3 \\
  &=\dfrac{1}{2}\Bigl(\dfrac{4+\dfrac{1}{a^2}}{2+\dfrac{\sqrt{3}}{a}}\Bigr)^3
\end{align}
である．
$a\to\infty$のとき括弧内は$\dfrac{4}{2}=2$に収束するから求める極限は
\begin{align}
\lim_{a\to\infty}\dfrac{T(a)}{S(a)}=\dfrac{1}{2}\cdot2^3=4
\end{align}
である．

{\bf [別解]}

\paragraph{(1)}
回転が$\dfrac{\pi}{6}$なので，$\tan$の加法定理を使って傾きを計算する方法も考えられる．
元の接線の傾きを$\tan\beta=2a$，$\ell$の接線の傾きを$\tan\alpha$とすると
\begin{align}
 \alpha = \beta - \dfrac{\pi}{6}
\end{align}
だから，
\begin{align}
 \tan\alpha 
   &= \dfrac{\tan\beta - \tan\dfrac{\pi}{6} }{1+\tan\beta\tan\dfrac{\pi}{6}} \\
   &= \dfrac{2a-\dfrac{\sqrt{3}}{3}}{1+\dfrac{2\sqrt{3}a}{3}} \\
   &= \dfrac{2\sqrt{3}a-1}{2a+\sqrt{3}}
\end{align}
となり，求める直線は$(a,a^2)$を通るから
\begin{align}
 y = \dfrac{2\sqrt{3}a-1}{2a+\sqrt{3}}\left(x-a\right)+a^2
\end{align}
である．



\newpage

\end{document}
